\section{Introduction}
\label{sec:introduction}

For a smooth projective variety $X$, canonical nonvanishing asks whether
pseudo-effectivity of $K_X$ guarantees a pluricanonical form: a nonzero
section of $mK_X$ for some integer $m>0$. Pseudo-effectivity says that the
numerical class is a limit of effective divisor classes. The question
concerns the canonical line bundle itself, rather than an approximation
to its numerical class. Over $\C$, pseudo-effectivity of $K_X$ is
equivalent to $X$ not being covered by rational curves
\cite[Corollary~0.3]{BDPP2013}. Nonvanishing thus asks for a
pluricanonical form on every smooth projective variety outside the
uniruled class.

Nonvanishing is weaker than abundance. For a nef canonical divisor,
abundance asks for \emph{semiampleness}: some positive multiple is
generated by its global sections and therefore defines a morphism.
Obtaining the first section is a separate difficulty. We prove the
following nonvanishing theorem in dimension four.

\begin{theorem}\label{thm:nonvanishing}
Let $X$ be a smooth connected projective complex fourfold. If $K_X$ is
pseudo-effective, then $H^0(X,mK_X)\ne0$ for some positive integer $m$.
\end{theorem}

The theorem imposes no condition on numerical dimension, irregularity,
or holomorphic Euler characteristic. Its conclusion is qualitative;
it gives no uniform bound on the integer $m$.

The threefold results explain both the history of the problem and the
lower-dimensional inputs used here. Miyaoka proved nonvanishing for
minimal threefolds \cite{Miyaoka1988Nonvanishing} and then abundance
in numerical dimension one \cite[pp.~203--204]{Miyaoka1988NuOne}.
Kawamata's work on pluricanonical systems and his abundance theorem for
minimal threefolds established the broader canonical case
\cite{Kawamata1985,Kawamata1992}. Keel--Matsuki--McKernan proved log
abundance for threefolds, with the correction included here
\cite{KMM1994,KMM2004}. Boundary arguments also require sections to
agree where components meet: Fujino's semi-log-canonical threefold
abundance theorem supplies generation on the whole reduced floor,
including its conductor identifications
\cite[Corollary~4.10]{Fujino2000}. These are established
lower-dimensional theorems used before the new fourfold argument.

Several important fourfold cases were already known. Fujino proved
semiampleness for canonical fourfolds with nef canonical divisor and
positive irregularity. He distinguished nonvanishing for non-uniruled
fourfolds of irregularity zero from the remaining abundance question
in Kodaira dimension zero
\cite[Corollary~4.7 and Problems~4.10--4.11]{Fujino2010}.
Lazi\'c--Peternell proved nonvanishing for $\Q$-factorial terminal projective minimal
varieties with canonical numerical dimension one and nonzero
holomorphic Euler characteristic
\cite[Theorem~6.7]{LazicPeternell2018}.
Ambro's canonical bundle formula gives a complementary reduction:
the lower-dimensional minimal model program and abundance settle a
nef klt adjoint whose nef reduction has lower-dimensional image
\cite[Theorem~4.3]{Ambro2005}. The nef reduction contracts the curves
of degree zero through very general points; the dimension of its image
is the \emph{nef dimension}. This curve-theoretic dimension is distinct
from numerical dimension, which is defined by nonzero intersection
powers of the nef divisor.

The smooth theorem also has a log canonical consequence. Hashizume
proved that smooth canonical nonvanishing in dimension $n$ implies
nonvanishing for log canonical pairs with real boundaries, as well as
log minimal models, in dimensions at most $n$
\cite[Theorem~1.4]{Hashizume2018}. For rational data we replace the
resulting effective real representative by a rational one, preserving
its support and working on the original variety.

\begin{corollary}\label[corollary]{cor:lc-nonvanishing}
Let $(X,\Delta)$ be a connected normal projective log canonical complex
pair of dimension at most four. Suppose that $\Delta\geq0$ is rational
and $D=K_X+\Delta$ is $\mathbb Q$-Cartier and nef. For every positive
integer $r$ such that $rD$ is Cartier, there is a positive integer $m$
with $H^0(X,\mathcal O_X(mrD))\ne0$.
\end{corollary}

The corollary does not require $X$ to be $\mathbb Q$-factorial, and
$m$ may depend on the pair and the prescribed index. Log abundance
asks for semiampleness of the nef adjoint $K_X+\Delta$. The passage
from the first section to that conclusion uses the supported-boundary
lifting and abundance-after-nonvanishing theorems in the separate
companion \emph{Lifting sections from the reduced support of an adjoint}
\cite[Theorems~1.1--1.2]{OpenAIBoundary2026}.
Together with the corollary, they give fourfold log abundance; that
consequence does not enter the proof below.

A recent comparison is Liu--Xu's preprint on good minimal models for
log canonical pairs of dimension at most five with nonnegative
invariant Iitaka dimension and numerical dimension at most one
\cite[Theorem~5.1]{LiuXu2025}. Its hypotheses already include
nonvanishing and restrict numerical dimension. The present smooth
theorem obtains the first section without either restriction.

\subsection*{Geometry of a possible counterexample}

Suppose nonvanishing fails. Existence of minimal models in dimension
four allows us to fix a terminal minimal model $Y$, retaining the
pluricanonical section spaces. The termination input used here is the
ordinary-pair case of Chen--Tsakanikas
\cite[Theorem 1.1]{ChenTsakanikas2023}; it supplies a minimal model,
without asserting semiampleness. Write $K=K_Y$. The divisor $K$ is nef
and not big. The preceding reductions give $q(Y)=h^1(Y,\mathcal O_Y)=0$
and full nef dimension. On a fixed very general locus, every curve has
positive $K$-degree and every proper positive-dimensional subvariety is
of general type. A fixed Cartier index makes the positive curve degrees
uniformly bounded away from zero.

Choose ample polarizations approaching the nef ray of $K$, scaled so
that their volumes remain small while their degrees on curves through
that locus tend to infinity. A section with unusually high vanishing
at a moving point would force a persistent component in the base locus
of the corresponding jet systems. Differentiation with respect to the
moving point controls the multiplicity along that component. This uses
the parameter-differentiation method of Ein--K\"uchle--Lazarsfeld
\cite[Proposition 2.3]{EinKuchleLazarsfeld1995}. We apply it to the
whole space of sections satisfying the jet conditions, so the estimate
controls ordinary powers of the ideal of the component. Intersection
theory bounds its degree and canonical intersection. General type then
produces a curve of bounded degree, contradicting the chosen scaling.
The resulting upper bound applies to every section in every allowed
Cartier degree.
The independent moving-center estimates and the finite-cover lemma
used in this argument and its two-factor extension are
\cite[Lemmas~7.1--7.5]{OpenAIAbundance2026}. We apply them to the
fourfold geometry and prove the required curve and correspondence
obstructions here.

A second use of this argument concerns sections over two copies of $Y$.
Fix an integer $\iota>0$ making $S_0=\iota K$ Cartier and consider
\[
 Z=\mathbb P\bigl(\mathcal O_{Y\times Y}(S_{0,1})
             \oplus\mathcal O_{Y\times Y}(S_{0,2})\bigr)
       \longrightarrow Y\times Y.
\]
A subscript indicates pullback from the corresponding factor. The two
summands give two distinguished sections of this projective bundle.
The symmetric-power formula records the distributions of a fixed
tensor exponent between the two factors: its summands are
$\mathcal O(aS_{0,1}+bS_{0,2})$, for nonnegative integers $a,b$
with $a+b$ fixed.
A polarization on $Z$ combines the two base polarizations and a multiple
of its tautological bundle. We show that it generates many jets at a
very general point away from the two distinguished sections. Its
Seshadri constant, the infimum of degree divided by multiplicity over
curves through the point, measures this local positivity.

Moving base components in $Z$ have several possible images and fibers.
One case requires an additional argument. Restrict to a slice where one
base coordinate is fixed. A multisection different from the two
distinguished sections meets them in divisors whose difference represents
a multiple of $K$. This is a \emph{signed representative}: its coefficients
may be positive or negative. We must prove that the logarithmic adjoint
on a resolution of its full support is big. This is where minimal metrics
and ordinary cohomological injectivity enter the geometry.

The companion \emph{Minimal metrics and interior injectivity
for nef adjoints}~\cite{OpenAIMetrics2026} supplies two analytic inputs. Every minimal semipositive metric on
the pullback of a nef projective klt adjoint has zero Lelong numbers.
In addition, an injectivity theorem on ordinary $H^1$ compares an
interior rational boundary with an endpoint boundary when the two endpoint
bundles carry such metrics. The exact hypotheses are recorded in
Section~\ref{sec:analytic-inputs}. These statements concern actual rational
line bundles; numerical equivalence would not suffice.

The cohomological method draws on the harmonic-form approach of Enoki,
as extended by Fujino to singular metrics, and on Matsumura's treatment
of metrics with transcendental singularities
\cite[Theorem~1.2, Lemmas~3.1--3.2, and Section~3, Claim~1]{FujinoInjectivity2012}
\cite[Theorems~1.3 and~5.9, Sections~5.2--5.3]{Matsumura2018}.
Fujino's complete metrics on an analytic complement and his comparison
with coherent cohomology with multiplier ideals give a model for the
setup. Matsumura supplies uniform estimates for primitives of exact forms.

Applied to two nearby nef klt adjoints, the cohomological injection
makes restriction to a whole reduced non-klt boundary surjective.
Lower-dimensional abundance and gluing across the conductor supply a
section on that boundary. Once this section has been lifted,
Gongyo--Matsumura's fourfold criterion converts nonvanishing and the
zero-Lelong metric into semiampleness
\cite[Corollary 5.3]{GongyoMatsumura2017}. This order matters: the
criterion is applied after the first section has been obtained.
The signed-representative argument then excludes the offending
multisection. The remaining analysis also excludes dominant finite
correspondences, using a fixed branch complement and the positive
canonical degree of curves.

\subsection*{From jets to a rank contradiction}

We now have opposing estimates. Sections on one copy of $Y$ can vanish
only to a small order, while the projective bundle over $Y\times Y$
has enough jets to give a large evaluation rank. We fix the polarization,
a smooth resolution and a flag of general hyperplane sections ending in
a curve, together with one finite system of jets. The scalar estimate
also supplies a fixed movable curve on the blowup of a point. Its
intersection with effective divisors will bound orders of sections
after reduction. We spread all these finite data and reduce modulo
sufficiently large primes.

Ordinary jets produce an evaluation matrix whose generic rank is large.
On the diagonal, the canonical Frobenius filtration expresses its
possible values through truncated powers of the cotangent bundle.
A single flag estimate bounds the sum of the dimensions of these spaces,
forcing a much smaller rank there. A nonzero maximal determinant must
therefore vanish to high order at a suitable diagonal point. Its two actual
determinant line bundles, one from each factor, have classes controlled
by the fixed movable curve. This bounds the order of every section of
either factor on the reduction itself. The resulting small determinant
order is incompatible with the order forced by the rank drop.

The filtration has a substantial history. For curves, the diagonal
filtration appears in Raynaud's work
\cite[Remarques~4.1.2(2)]{Raynaud1982}; Joshi--Ramanan--Xia--Yu
formulate it using the Cartier connection
\cite[Section 5.3]{JoshiRamananXiaYu2006}. Sun and Kitadai--Sumihiro
give the higher-dimensional filtration and its local description by
truncated monomials \cite[Theorem 3.7]{Sun2008}
\cite[Corollary 3.5]{KitadaiSumihiro2008}.
Musta\c{t}\u{a}--Schwede's comparison of ordinary and Frobenius powers
provides the local jet conversion \cite[proof of Proposition~2.12, Equation~(2.8)]{MustataSchwede2014},
and Langer's estimate controls the Frobenius instability error
\cite[Corollary 2.5]{Langer2004}. The argument combines these tools
with the two-factor geometry and the common flag estimate.
Section~\ref{sec:frobenius} specifies the numerical separation between
the resulting ranks and orders and verifies the inputs of the common theorem.

The complete finite-data Frobenius comparison is proved independently in
\cite[Theorem~9.1]{OpenAIAbundance2026}. It
is also applicable to the all-dimensional abundance argument. Its hypotheses are the stated
geometric data and numerical inequalities. No all-dimensional abundance
conclusion or logarithmic Iitaka subadditivity is used in the fourfold
preparation of these data.

\paragraph{Reading order.}
Section~\ref{sec:reduction} fixes the terminal model and its curve-degree
obstruction. Section~\ref{sec:moving-centers} develops the moving-center
tools and proves the scalar order bound. Section~\ref{sec:analytic-inputs}
states the analytic inputs, and Section~\ref{sec:signed-representatives}
proves logarithmic general type for signed supports.
Section~\ref{sec:two-slot-jets} excludes the product's moving centers
and obtains the two-slot jets.
Section~\ref{sec:frobenius} verifies the finite data for the common
Frobenius theorem, completing smooth nonvanishing.
Section~\ref{sec:conclusion} gives the original-variety and
prescribed-index consequence.

\paragraph{Conventions.}
Varieties are over $\C$ except when the finite-data comparison is
reduced to positive characteristic. A subvariety is of general type
when a smooth projective resolution is of general type. Rational
line bundles and their sections are interpreted after clearing
specified denominators. An inequality in nef order means that its
difference is nef. A very general locus is the complement of a
countable union of proper closed subsets. The symbols $r$ for a
prescribed Cartier index, $r_t$ for a scaling integer, $q(Y)$ for
irregularity, and $q$ for a projective-bundle weight have distinct uses.
